% 06_reproducir_experimentos.tex - Manual Capítulo 6: Reproducir Experimentos
\chapter{Reproducir los Experimentos Reportados}
\label{man:reproducir}

Todos los resultados del informe técnico (\texttt{docs/results/}) se pueden regenerar
ejecutando los comandos exactos que se listan a continuación. Los valores de conteo
(generados, expandidos, frontera, costo, longitud) son \textbf{determinísticos} e idénticos
en cada ejecución; solo \code{tiempo\_segundos} varía (tiempo de ejecución).

\section{Preparación}
\label{sec:prep-reproducir}

\begin{minted}{bash}
cd tactical-grid-ai
uv sync                    # instala dependencias
uv run pytest              # verifica que 281 tests pasan
R=docs/results
U="uv run python -m tactical_grid"
mkdir -p $R
\end{minted}

\section{Casos 1--2: Búsqueda No Informada y Costos (F2)}
\label{sec:rep-f2}

\begin{minted}{bash}
# Caso 1: BFS vs DFS
$U scenarios/caso_01_bfs_vs_dfs.json --algorithm BFS,DFS,UCS --verify \
  --csv $R/caso_01_bfs_vs_dfs.csv --json $R/caso_01_bfs_vs_dfs.json

# Caso 2: BFS vs UCS (ruta corta vs económica)
$U scenarios/caso_02_bfs_vs_ucs.json --algorithm BFS,DFS,UCS --verify \
  --csv $R/caso_02_bfs_vs_ucs.csv --json $R/caso_02_bfs_vs_ucs.json
\end{minted}

\textbf{Valores esperados (derivados a mano, asertados en tests):}

\begin{table}[H]
\centering
\caption{Caso 1: BFS vs DFS}
\begin{tabular}{lccccccc}
\toprule
Alg & Costo & Movs & Gen & Exp & MaxF & Match \\
\midrule
BFS & 4 & 4 & 39 & 29 & 11 & Sí \\
DFS & 42 & 42 & 152 & 94 & 58 & Sí \\
UCS & 4 & 4 & 39 & 29 & 11 & Sí \\
\bottomrule
\end{tabular}
\end{table}

\begin{table}[H]
\centering
\caption{Caso 2: BFS vs UCS}
\begin{tabular}{lccccccc}
\toprule
Alg & Costo & Movs & Gen & Exp & MaxF & Match \\
\midrule
BFS & 25 & 7 & 126 & 104 & 23 & Sí \\
DFS & 153 & 87 & 396 & 287 & 171 & Sí \\
UCS & 17 & 17 & 269 & 239 & 33 & Sí \\
\bottomrule
\end{tabular}
\end{table}

\section{Casos 3--5: Búsqueda Informada y Beam (F3)}
\label{sec:rep-f3}

\begin{minted}{bash}
# Caso 3: Obstáculo + heurística
$U scenarios/caso_03_obstaculo_heuristica.json \
  --algorithm UCS,A_STAR --heuristic H1,H2 --verify \
  --csv $R/caso_03_obstaculo_heuristica.csv --json $R/caso_03_obstaculo_heuristica.json

# Caso 4: Comparación heurísticas
$U scenarios/caso_04_comparacion_heuristicas.json \
  --algorithm UCS,A_STAR --heuristic H1,H2 --verify \
  --csv $R/caso_04_comparacion_heuristicas.csv --json $R/caso_04_comparacion_heuristicas.json

# Caso 5: Beam k=1,2,4,8
$U scenarios/caso_05_beam_k.json \
  --algorithm UCS,BEAM --heuristic H1 --k 1,2,4,8 --verify \
  --csv $R/caso_05_beam_k.csv --json $R/caso_05_beam_k.json
\end{minted}

\textbf{Valores esperados:}

\begin{table}[H]
\centering
\caption{Caso 3: UCS, A* h1, A* h2}
\begin{tabular}{lcccccc}
\toprule
Alg & Costo & Movs & Gen & Exp & MaxF & Match \\
\midrule
UCS & 29 & 29 & 307 & 295 & 17 & Sí \\
A* h1 & 29 & 29 & 194 & 166 & 28 & Sí \\
A* h2 & 29 & 29 & 122 & 84 & 38 & Sí \\
\bottomrule
\end{tabular}
\end{table}

\begin{table}[H]
\centering
\caption{Caso 4: UCS, A* h1, A* h2}
\begin{tabular}{lcccccc}
\toprule
Alg & Costo & Movs & Gen & Exp & MaxF & Match \\
\midrule
UCS & 24 & 22 & 346 & 332 & 20 & Sí \\
A* h1 & 24 & 22 & 230 & 186 & 44 & Sí \\
A* h2 & 24 & 22 & 164 & 115 & 49 & Sí \\
\bottomrule
\end{tabular}
\end{table}

\begin{table}[H]
\centering
\caption{Caso 5: Beam k=1,2,4,8}
\begin{tabular}{cccccccc}
\toprule
k & Éxito & Costo & Movs & Gen & Exp & MaxF & Match \\
\midrule
1 & No & -- & -- & 9 & 5 & 3 & Sí \\
2 & Sí & 30 & 16 & 43 & 29 & 6 & Sí \\
4 & Sí & 28 & 14 & 50 & 41 & 6 & Sí \\
8 & Sí & 28 & 14 & 46 & 44 & 6 & Sí \\
\bottomrule
\end{tabular}
\end{table}

\section{Casos 6--10: Adversarial (F4)}
\label{sec:rep-f4}

\begin{minted}{bash}
# Caso 6: Estado compuesto (test unitario)
uv run pytest tests/test_case_6.py -v

# Caso 7: Minimax profundidades 1-3
$U scenarios/caso_07_decision_adversarial.json \
  --algorithm MINIMAX --depth 1,2,3 --verify \
  --csv $R/casos_07_10_adversarial.csv --json $R/casos_07_10_adversarial.json

# Caso 8: Minimax frente a alfa-beta, profundidad 4
$U scenarios/caso_07_decision_adversarial.json \
  --algorithm MINIMAX,ALPHA_BETA --depth 4 --verify \
  --csv $R/casos_07_10_adversarial.csv --json $R/casos_07_10_adversarial.json

# Caso 9: Orden de acciones (profundidad 3)
$U scenarios/caso_07_decision_adversarial.json \
  --algorithm ALPHA_BETA --depth 3 --order default,reversed --verify \
  --csv $R/casos_07_10_adversarial.csv --json $R/casos_07_10_adversarial.json

# Caso 10: Profundidades 1-5 (alfa-beta)
$U scenarios/caso_07_decision_adversarial.json \
  --algorithm ALPHA_BETA --depth 1,2,3,4,5 --verify \
  --csv $R/casos_07_10_adversarial.csv --json $R/casos_07_10_adversarial.json
\end{minted}

\textbf{Valores esperados (casos 7-10):}

\begin{table}[H]
\centering
\caption{Caso 7: Minimax profundidades 1-3}
\begin{tabular}{lccc}
\toprule
Prof & Acción & Valor & Match \\
\midrule
1 & A2$\to$(0,1) & 3 & Sí \\
2 & A2$\to$(0,1) & 0 & Sí \\
3 & A1$\to$(1,1) & 3 & Sí \\
\bottomrule
\end{tabular}
\end{table}

\begin{table}[H]
\centering
\caption{Caso 8: Minimax vs Alfa-beta (prof 4)}
\begin{tabular}{lcccccc}
\toprule
Alg & Acción & Valor & Gen & Eval & Podados & Match \\
\midrule
Minimax & A1$\to$(0,0) & 0 & 4826 & 4296 & 0 & Sí \\
Alpha-beta & A1$\to$(0,0) & 0 & 431 & 299 & 757 & Sí \\
\bottomrule
\end{tabular}
\end{table}

\begin{table}[H]
\centering
\caption{Caso 9: Orden default vs reversed (prof 3)}
\begin{tabular}{lcccccc}
\toprule
Orden & Acción & Valor & Gen & Eval & Podados & Match \\
\midrule
default & A1$\to$(1,1) & 3 & 142 & 110 & 106 & Sí \\
reversed & A1$\to$(1,1) & 3 & 191 & 148 & 145 & Sí \\
\bottomrule
\end{tabular}
\end{table}

\begin{table}[H]
\centering
\caption{Caso 10: Alfa-beta profundidades 1-5}
\begin{tabular}{lcccc}
\toprule
Prof & Acción & Valor & Eval & Match \\
\midrule
1 & A2$\to$(0,1) & 3 & 7 & Sí \\
2 & A2$\to$(0,1) & 0 & 24 & Sí \\
3 & A1$\to$(1,1) & 3 & 110 & Sí \\
4 & A1$\to$(0,0) & 0 & 299 & Sí \\
5 & A1$\to$(1,1) & 3 & 1583 & Sí \\
\bottomrule
\end{tabular}
\end{table}
(Minimax evals: 7 / 64 / 467 / 4296 / 30994)

\section{Rendimiento (G13)}
\label{sec:rep-rendimiento}

\begin{minted}{bash}
# 15x15 mapa (3 unidades/bando)
$U scenarios/prueba_15x15.json --algorithm ALPHA_BETA --depth 5,6,7 \
  --csv $R/rendimiento_15x15.csv

$U scenarios/prueba_15x15.json --algorithm MINIMAX --depth 4,5,6 \
  --csv $R/rendimiento_15x15_minimax.csv

# 20x20 mapa
$U scenarios/prueba_20x20.json --algorithm ALPHA_BETA --depth 2,3,4,5,6 \
  --csv $R/rendimiento_20x20.csv
\end{minted}

Valores esperados (ver la sección ``Rendimiento (G13)'' del informe técnico).


\section{Sweep Completo (Todos los Casos 1-10 + Apéndice A)}
\label{sec:rep-sweep}

\begin{minted}{bash}
R=docs/results; U="uv run python -m tactical_grid"
$U scenarios/apendice_a.json --algorithm BFS,DFS,UCS,A_STAR,BEAM \
  --heuristic H1,H2 --k 1,2,4,8 --verify \
  --csv $R/apendice_a.csv --json $R/apendice_a.json
\end{minted}

\section{Caso 11: Variante dinámica}
\label{sec:rep-dinamica}

\begin{minted}{bash}
# Comparar Minimax dinámico y Alfa-beta dinámico
$U scenarios/caso_11_dinamico.json \
  --variant dinamica --algorithm MINIMAX,ALPHA_BETA --depth 1,2,3,4 --verify \
  --csv $R/caso_11_dinamico.csv --json $R/caso_11_dinamico.json

# Ejecutar una partida completa con la variante dinámica
$U scenarios/caso_11_dinamico.json \
  --mode partida --play auto --variant dinamica --algorithm ALPHA_BETA --depth 3
\end{minted}

En la variante dinámica se reportan además los efectos de recuperación e inmunidad del
portador. Las métricas de tiempo pueden cambiar entre ejecuciones; las acciones y conteos
son deterministas para el mismo escenario y parámetros.

\section{Verificación de que Todo Pasa}
\label{sec:verificar-todo}

\begin{minted}{bash}
# Tests completos
uv run pytest                    # 281 tests, 0 skips
uv run ruff check .              # lint
uv run ruff format --check .     # formato
uv run mypy src                  # tipos strict

# Verificar archivos de resultados generados
ls -la docs/results/
# Debe existir: caso_01_bfs_vs_dfs.csv/json, caso_02_bfs_vs_ucs.csv/json, ...
# caso_03_obstaculo_heuristica.csv/json, caso_04_comparacion_heuristicas.csv/json
# caso_05_beam_k.csv/json, casos_07_10_adversarial.csv/json
# caso_11_dinamico.csv/json
# rendimiento_15x15.csv, rendimiento_20x20.csv, apendice_a.csv/json
\end{minted}

Todos los valores de conteo (generados, expandidos, frontera, costo, longitud) deben
coincidir \textbf{exactamente} con las tablas de este manual y del informe técnico
(\texttt{cases.md}). Solo \code{tiempo\_segundos} varía entre ejecuciones.